\documentclass[quick,con]{debate}
\motion{「勇敢的人先享受世界」是一个陷阱 / 不是一个陷阱}
\stance{反方}{不是一个陷阱}
\begin{document}
\debatetitle
\begin{thesisbox}它给的是一个顺序，不是一张支票；说好的「更早开始」兑现了，没兑现的那部分是准备的事，不是这九个字的事。\end{thesisbox}

\section{定义与判准（标准表述）}
\begin{fields}
\field{陷阱}{它呈现出来的样子和照做之后的真实后果系统性地不一致，而且正是这层外观把人引了进去。\sub{一辩必说}{我方承认陷阱可以没有设局者。绝不说「必须有人设局」，那是定义杀。}}
\field{勇敢的人}{在后果不确定时仍然选择行动的人；勇气针对的是已知的风险，不知道有风险还冲的叫莽撞。}
\field{它说了什么（正面回答）}{在先动和先等之间，它推荐先动。不是「它什么都没说」。}
\field{判准}{三件事：① 外观与真实因果系统性不一致；② 照做的人受损是常态；③ 受损正是被它遮住的东西造成的。\sub{我方要证明}{打掉任意一件即可，今天打第一件和第三件。}\sub{对方要证明}{三件都成立，举证责任全在对方。}}
\field{对方提偏向定义或判准时的 30 秒口径}{对方若说「不必统计受损」：那「努力就有回报」也是陷阱，这条标准量的是有没有例外。我方承认它不严谨，不严谨和圈套隔着三件事。}
\end{fields}

\section{我方论点}
\begin{tabularx}{\linewidth}{@{}L{4.6em}Y Y L{7em}@{}}\toprule
\thead{标签}&\thead{一句话}&\thead{核心数据 / 学理 / 案例}&\thead{出处}\\\midrule
\textbf{顺序不是支票}&它承诺的只有「先行动的人更早开始」，这件事兑现了，没有落差&数据：34.8\% 有跳槽意愿者会裸辞，其中 21.2\% 是「休息后再出发」，21.8\% 会 GAP；学理：陷阱要么有骗，要么有出不来的坑底；案例：最早是旅游九宫格配文&智联招聘 2025 春季跳槽调研，37980 份样本；汉典「陷阱」词条；澎湃 2023-10-17\\
\textbf{账记在准备上}&让人摔的是没留退路，不是有勇气；这九个字没有一个字让你别算账&数据：长期回顾 51\% 困扰于做过的事，49\% 困扰于没做的事；学理：亲历的成功是能力感最可靠的来源&Richardson \& Gilovich 2023, n=2600；Bandura 1977\\
\bottomrule\end{tabularx}

\section{对方论点 → 一句话反驳}
\begin{tabularx}{\linewidth}{@{}Y Y Y@{}}\toprule
\thead{对方会说}&\thead{我方一句话回}&\thead{对方追问时}\\\midrule
门票换成了勋章&你方自己的数据里，四成多人清清楚楚算了账；省略是把话说短，盖住是把东西藏起来&说「他们被判成不勇敢」→ 判他们的是你方的读法，原句没说不勇敢的人是谁\\
删过的账本、幸存者偏差&那学长的经验帖、课本上的科学家故事也全是陷阱&说「区别在顶替」→ 请证明有人真因为这句话放弃算账\\
青年失业率高，空窗代价大&代价高说明要准备，不说明它是坑；而且这个数还在回落&说「对没兜底的人不可逆」→ 承认，所以要准备，这正是我方第二点\\
说到底这就是消费主义话术&贴标签不等于证明性质，三件事一件都没回答&说「消费主义就是伪装」→ 那请说出它向谁收了钱、承诺了什么\\
\bottomrule\end{tabularx}

\section{必问与必答}
\begin{points}{必问（对方不答就点名、重复、不换题）}
\item 这句话盖住的是哪几个字？（主问题）
\item 「受损是常态」，数据在哪？
\item 省略和顶替，除了你方的读法还有什么分得开？
\end{points}
\begin{points}{必答（15 秒正面回答）}
\item 「享受两个字你方删哪去了？」→ 没删。享受是结果的名字，不是保证；原句没说享受得到多少、留多久。
\item 「一句什么都没说的话怎么会红三年？」→ 我方没说它什么都没说，它说的是别把一生花在等条件齐。
\item 「你方承不承认有人因此摔了？」→ 承认。摔的是没留退路的人，那是准备的事。
\end{points}

\section{自由辩战场概要}
\begin{tabularx}{\linewidth}{@{}L{4.2em}Y Y Y@{}}\toprule
\thead{战场}&\thead{开场问题}&\thead{目标}&\thead{收束语}\\\midrule
伪装（前 90 秒）&盖住的是哪几个字？&对方始终交不出被盖住的东西&空的一格撑不起一个陷阱\\
常态（中 90 秒）&照它做的人变差，数据在哪？&对方自立的第二件事全场无证据&三件缺一，它就只是一句不严谨的话\\
\bottomrule\end{tabularx}
时间：前 90 秒战场一并抛完必问，中 90 秒战场二，后 60 秒只重复主问题、不开新战场。站位：一辩守定义判准，二辩接驳论，三辩接盘问，四辩收束。

\textbf{各环节}：1 反一立论 180 秒，开头先让出「陷阱不需要设局者」；3 反四质询正一 120 秒双边，逼对方交出被盖住的字；4 反二驳论 120 秒，先驳，打对方最难圆的第二件事；6 反三盘问 90 秒，先问正二再问正四同一问题；8 反三小结 90 秒，留临场位回应正三；10 反四结辩 210 秒，封路式，先拆掉正四的故事。被质询与被盘问时每答不超过八秒。

\section{如果……就……}
\begin{contingency}
\ifthen{对方说我方在替这句话辩护}{立刻回：我方一句都没夸过它，我方只说它不满足陷阱的要件。}
\ifthen{队友说漏了「必须有设局者」}{下一次发言第一句更正：那一格我方一辩就让了。}
\ifthen{「顺序不是支票」被打穿}{收缩到「账记在准备上」，用判准第三件解释为什么仍然占优。}
\ifthen{对方回避必问或拿出没预料到的数据}{回避就点名、重复、不换题；新数据先问出处与自变量。}
\end{contingency}
\end{document}
